\section{Introduction}
Uplift modeling refers to the set of techniques used to estimate the effect of an action on a user's outcome. This technology can be applied to various fields, such as economics\cite{olaya2020survey,gubela2017revenue}, medicine\cite{jaroszewicz2014uplift,jaskowski2012uplift} and sociology\cite{olaya2020uplift,guo2020counterfactual}. For example, an e-commerce company is preparing to send promotional coupons (i.e., action or treatment) to some users to attract them to purchase more products (i.e., outcome). Then estimating the individual uplift can help to find target users. Uplift modeling is a complicated problem because one needs to estimate the difference between two outcomes with and without the treatment which are mutually exclusive to individuals. For example, we can only observe the outcome of a user getting or not getting a promotional coupon. The outcome we can observe is called the factual outcome and the outcome we can not observe is called the counterfactual outcome. Therefore, uplift modeling also can be seen as a counterfactual inference problem\cite{gutierrez2017causal,kuang2020data}. 

To model the counterfactual outcome, existing uplift methods mainly rely on randomized experiments or observational data: users will be assigned to either the treatment group or the control group and we can observe only one type of outcome. Based on these data, most existing methods \cite{gutierrez2017causal,guelman2015uplift,rzepakowski2012decision} utilize user's individual features to estimate the user uplift \cite{breiman2001random,chen2016xgboost}. An accurate uplift modeling is required to capture the complex and the unobserved hidden factors (representations) within the user. However, many factors regarding the user's uplift are difficult to be captured by only using individual data for two reasons. Firstly, in real-world scenarios, the individual features of users, especially the new users, are missing. Secondly, some informative information for uplift estimation is hidden and difficult to be characterized by individual features, like social status and personality.

To resolve the aforementioned problems, we hope to introduce the social graph. On the one hand, considering that users with close social relations usually have similar behaviors and preferences, we can utilize the information from social neighbors as a supplement to the user's own features. On the other hand, social graphs can reveal some informative social information like social status, which is also beneficial for uplift modeling \cite{veitch2019using}. From the experiments, we also demonstrate that the uplift difference between users with friend relationships is much smaller than the difference between random users. Therefore, it is very essential to incorporate social graphs when doing uplift estimation. 

A few of the existing works propose graph-based uplift models for uplift estimation \cite{veitch2019using,guo2020counterfactual,guo2020learning}: They first use graph neural network (GNN) model to learn the graph-based representations from the social graphs. Then with the graph-based representations, these works design uplift estimators, which consist of two separate pathways of models to predict the outcome with and without the treatment using the data of the treatment group and control group respectively. Although these works have achieved substantial improvements, they still have limitations. In real-world scenarios, since the effect of the treatment is unpredictable, imposing the treatment on samples may bring a bad effect like capital loss and user loss. In this way, it is common that only a little flow will be assigned to the treatment group. In this way,  the number of labeled instances, particularly the instances from the treatment group, is commonly limited, especially in the scenario of randomized experiments we mainly focus on. Moreover, to model the relational data, graph-based models require learning more parameters. From Figure \ref{sec:scarcity}, compared with tree-based and NN-based models, the performance of existing graph-based models will drop more quickly when the labeled data is scarce, thereby labeled data scarcity is a more severe and unresolved problem for graph-based representation learning in uplift estimation. 

To address aforementioned problems, we propose a general GNN-based framework with two uplift estimators for uplift modeling. First of all, a GNN-based model with breadth and depth aggregators is proposed to generate the graph-based representations for the following uplift estimation. And empirically we demonstrate the proposed uplift estimation framework is general for many existing GNN-based models. Furthermore, to address the label scarcity problem, we propose two uplift estimators to utilize the treatment and control group data and their corresponding social relations. Firstly, we design a class-transformed target, which we prove is equal to the uplift and is general for uplift scenarios. Unlike two separate pathways of estimators, our transformed target is able to utilize the training instances with and without the treatment simultaneously in a common model to approach the uplift. Furthermore, when the outcome is discrete, we design partial labels based on the user's treatment and the observed outcome. Then we design the other uplift estimator with two classifiers to learn partial labels. The two classifiers can focus on different facets of the uplift but all require two groups of data for training. In this way, our model can capture the relations between two groups of data and utilize the labeled data more effectively. Experimental results show our proposed framework can outperform the best-performing baseline method by an improvement of $5$\% to $10$\% in the regression setting and $12$\% to $25$\% in the classification setting. 

The main contributions of the paper are summarized here:
\begin{itemize}
    \item \textbf{Problem}: We point out the label scarcity problem is severe for uplift modeling, especially when trained with the graph-based model. To the best of our knowledge, it is the first work trying to solve the label scarcity problem in uplift estimation.
    \item \textbf{Methodology}: We introduce a novel GNN-based framework with two uplift estimators for uplift modeling, which can utilize the two groups of data and the social graph comprehensively. Specifically, when the outcome is discrete, we are the first to introduce partial label learning to uplift estimation, which is able to utilize the labeled data more effectively to alleviate the labeled data scarcity problem.
    \item \textbf{Results}: Extensive experiments on a public dataset and two industrial datasets demonstrate the superiority of proposed GNN-based uplift model(GNUM) on different types of outcomes. Specifically, we find labeled data scarcity is indeed a serious problem for previous graph-based uplift methods and our methods are robust to the label scarcity problem. 
\end{itemize}
